\documentclass[10pt,a4paper]{amsart}
\usepackage[margin=19mm]{geometry}
\usepackage{amsmath,amssymb,mathtools}
\usepackage[scr=rsfs,cal=euler]{mathalfa}
\usepackage{xfrac}
\usepackage[unicode,hidelinks,bookmarks=false]{hyperref}
\newcommand{\ie}{i.e.\ }
\newcommand{\cf}{cf.\ }
\newcommand{\resp}{respectively\ }
\newcommand{\Subsection}{\subsection}
\providecommand{\nobreakdash}{\nobreak\hbox{-}}
\setlength{\emergencystretch}{2em}
\multlinegap0pt
\usepackage{bm}
\usepackage{bbm}
\usepackage{dsfont}
\usepackage{xspace}
\usepackage{cancel}
\usepackage{mathtools}
\usepackage{cleveref}
\usepackage{enumitem}
\usepackage{amsthm}
\makeatletter
\@ifundefined{equation}{\newtheorem{equation}{Equation}}{}
\@ifundefined{definition}{\newtheorem{definition}[equation]{Definition}}{}
\@ifundefined{proposition}{\newtheorem{proposition}[equation]{Proposition}}{}
\@ifundefined{remark}{\newtheorem{remark}[equation]{Remark}}{}
\makeatother
\newcommand{\Cat}{\mathscr{C}\mathrm{at}}
\newcommand{\Fun}{\mathrm{Fun}}
\newcommand{\Hom}{\mathrm{Hom}}
\newcommand{\Mor}{\mathrm{Mor}}
\newcommand{\Obj}{\mathrm{Obj}}
\newcommand{\cC}{\mathcal{C}}
\newcommand{\cD}{\mathcal{D}}
\newcommand{\id}{\mathrm{id}}
\newcommand{\sSet}{s\mathscr{S}\mathrm{et}}
\title[Introduction to Complete Segal Spaces]{Introduction to Complete Segal Spaces}
\author{Nima Rasekh}
\date{Source: arXiv:1805.03131v2 (CC BY 4.0)}
\begin{document}
\maketitle

\noindent\emph{Source and licence.} Introduction to Complete Segal Spaces. Version arXiv:1805.03131v2 (CC BY 4.0), distributed under the Creative Commons Attribution 4.0 International licence, as recorded in the arXiv licence metadata for this exact version and shown on the abstract page https://arxiv.org/abs/1805.03131v2. Full rights notice: licenses/arxiv-1805.03131-CC-BY-4.0.txt.\par\smallskip

\noindent\textit{Editorial scope.} Counted unit: the Proposition whose source label is \texttt{prop:nerve fully faithful} in the article (source0.tex lines 1122--1124, source label \texttt{prop:nerve fully faithful}), delivered together with the article's own complete original proof of it (source0.tex lines 1126--1164, one \texttt{proof} environment of 39 lines). No mathematics is written, repaired or re-derived: statement and proof are the article's own text line for line. Required context retained uncounted: rem:nerve explicit (lines 1101--1110); def:fully faithful (lines 4180--4186). Every definition, display and label of those lines is retained unchanged. Omitted from this package and not redistributed: 1--1100, 1111--1121, 1125--1125, 1165--4179, 4187--4362. Nothing inside a retained excerpt is omitted or altered: every retained body is delivered exactly as the article's lines are written, so no omission receipt of any kind accompanies this package. Citations: the retained excerpts carry no citation command at all, so no bibliography entry of the article is transcribed and no reference is redistributed. Reference adaptation: none was needed and none was made; every reference inside the delivered unit resolves inside it, so no unresolved reference or citation is produced. Printed slips kept literal: no printed word, symbol, hypothesis or number of the counted statement or of its proof was altered. Layout and class: the article's own class is \texttt{amsart} with packages including amssymb, amsfonts, amsmath, amsthm, xy, enumerate, mathrsfs, mathtools, euscript, array, enumitem, tikz, tikz-cd, textcomp; the wrapper is the lane's pinned \texttt{amsart} editorial wrapper at 10pt on A4, which restates the theorem-like environments the retained text needs and adds the article's own macro definitions that the retained text needs (\texttt{\textbackslash Cat}, \texttt{\textbackslash Fun}, \texttt{\textbackslash Hom}, \texttt{\textbackslash Mor}, \texttt{\textbackslash Obj}, \texttt{\textbackslash cC}, \texttt{\textbackslash cD}, \texttt{\textbackslash id}, \texttt{\textbackslash sSet}), with extra wrapper packages (bm, bbm, dsfont, xspace, cancel, mathtools, cleveref, enumitem). The delivered numbering is the unit's own. Assets: no retained span contains a figure environment, a graphics-inclusion command, a PDF-page inclusion, a table or a TikZ picture; no image file is copied into this package, the package declares no asset dependency and no image file is redistributed with it. Licence: version arXiv:1805.03131v2 of this article is distributed under the Creative Commons Attribution 4.0 International licence, as recorded in the arXiv licence metadata for this exact version and shown on the abstract page https://arxiv.org/abs/1805.03131v2. A full rights notice is delivered at licenses/arxiv-1805.03131-CC-BY-4.0.txt. Characters: every retained excerpt is pure ASCII, so the delivered engine, XeLaTeX with its default Latin Modern text font, reports no missing character.
\begin{remark} \label{rem:nerve explicit}
  More explicitly, we can describe the nerve as follows. 
  \begin{itemize}
    \item $n = 0$: $N\cC_0 = \Fun([0],\cC) = \Obj_{\cC}$ is the set of objects of $\cC$.
    \item $n = 1$: $N\cC_1 = \Fun([1],\cC) = \Mor_{\cC}$ is the set of morphisms of $\cC$.
    \item $n \geq 2$: $N\cC_n = \Fun([n],\cC)$ is the set of $n$ composable morphisms in $\cC$. In other words, an element of $N\cC_n$ is a sequence of morphisms
    \[x_0 \xrightarrow{f_1} x_1 \xrightarrow{f_2} x_2 \xrightarrow{f_3} ... \xrightarrow{f_n} x_n\]
    where $x_i \in \Obj_\cC$ and $f_i\colon x_{i-1} \to x_i$.
  \end{itemize}
\end{remark}

\begin{definition} \label{def:fully faithful}
  A functor $F\colon \cC \to \cD$ is called \emph{fully faithful} if for every pair of objects $X,Y$ in $\Obj_{\cC}$, the map 
  \[
  F_{X,Y}\colon \Hom_{\cC}(X,Y) \to \Hom_{\cD}(F(X),F(Y))
  \]
  is a bijection of sets.
\end{definition}

\begin{proposition} \label{prop:nerve fully faithful}
 The nerve is fully faithful (\cref{def:fully faithful}).
\end{proposition}

\begin{proof}
  Let $\cC$ and $\cD$ be two categories. We need to show that the map
  \[
  N\colon \Hom_{\Cat}(\cC,\cD) \to \Hom_{\sSet}(N\cC,N\cD)
  \]
  is a bijection. Observe that for a given functor $F\colon \cC\to \cD$, the map $N(F)_0\colon \Obj_{\cC} \to \Obj_{\cD}$ is the functor $F$ on objects, and $N(F)_1\colon \Mor_{\cC} \to \Mor_{\cD}$ is the functor $F$ on morphisms. Hence, for two functors $F,G\colon \cC\to \cD$ with $N(F)=N(G)$, we have $F_0=G_0$ and $F_1=G_1$, meaning $F = G$. Hence $N$ is injective.

  It remains to show that $N$ is surjective. Given a map $f\colon N\cC \to N\cD$ of simplicial sets, we will define a functor $F\colon \cC\to \cD$, such that $N(F)=f$. We define $F$ as follows.
  \begin{itemize}[leftmargin=*]
    \item For an object $x\in \Obj_{\cC}=(N\cC)_0$, define $F(x)\coloneq f_0(x)$.
    \item For a morphism $g\colon x\to y$ in $\cC$, define $F(g) \coloneq f_1(g)$.
    \item[] Here the simplicial relations
    \[
      0^*(f_1(g))=f_0(0^*(g))=f_0(x)=F(x),\qquad
      1^*(f_1(g))=f_0(1^*(g))=f_0(y)=F(y),
    \]
    imply that if $g\colon x\to y$, then $F(g)\colon F(x)\to F(y)$.
    \item For a given object $x$ in $\cC$, 
    \[
      F(\id_x) =  f_1(00^*(x)) = 00^*(f_0(x)) = 00^*(F(x)) = \id_{F(x)},
    \]
    hence $F$ preserves identities.
    \item For composable morphisms $x\xrightarrow{g}y\xrightarrow{h}z$ in $\cC$, we have
    {\small
    \[
      F(h\circ g) = f_1(h\circ g) = f_1(02^*(g,h)) = 02^*(f_2(g,h)) = 02^*(f_1(g),f_1(h)) = f_1(h) \circ f_1(g) = F(h)\circ F(g),
    \]
    }
    hence $F$ preserves composition.
  \end{itemize}

  Finally, we confirm that $N(F)=f$. By definition $N(F)_0 = f_0$ and $N(F)_1 = f_1$. More generally let $\sigma\in (N\cC)_n$ be an $n$-simplex, which, by \cref{rem:nerve explicit}, is a string of $n$ composable morphisms
  \[
  \sigma=(x_0\xrightarrow{g_1}x_1\xrightarrow{g_2}\cdots\xrightarrow{g_n}x_n).
  \]
  Then the simplicial identities imply that 
  \[ f_n(\sigma) = (f_1(g_1),..., f_1(g_n)) = (F(g_1),..., F(g_n)) = N(F)_n(\sigma). \]
  This finishes the proof.
\end{proof}

\end{document}
